% MIA5150 course proposal. Edit this file; GitHub Actions rebuilds its PDF.
% Measured results come from skillret-pilot/reference_run/, not future experiments.
\documentclass{article}
\PassOptionsToPackage{numbers,compress}{natbib}
% Preprint shows team names without implying acceptance at NeurIPS.
\usepackage[preprint]{neurips_2026}
\usepackage[utf8]{inputenc}
\usepackage[T1]{fontenc}
\usepackage{microtype}
\usepackage{booktabs}
\usepackage{amsmath,amsfonts}
\usepackage{graphicx}
\usepackage[hidelinks]{hyperref}
\usepackage{url}

\title{Agent Skill Discovery at Scale}
\author{Research Team}

\begin{document}
\maketitle

\begin{abstract}
An agent with thousands of available skills must decide what to load without
spending its context budget on the whole library. We propose a controlled
comparison of listing, vector retrieval and domain-then-vector retrieval, with
BM25 as a lexical reference. Library size is the main experimental variable;
selection quality and model token use are the main outcomes. An executed pilot
established a full-document BM25 Recall@5 of 51.67\% on 200 queries over 6,006
skills, and verified native Deep Agents integration with five OpenRouter models.
The small live sample exposed invalid selections and provider limits rather than
a reliable model ranking. The next study will isolate discovery strategy under
fixed inputs, a pinned model/provider and an explicit request budget.
\end{abstract}

% Add student numbers here if the course requires them; none were supplied.
\noindent\textbf{MIA5150 project proposal --- Benchmark and Empirical Study.}
\textbf{Research claim to test:} skill shortlisting reduces model context cost,
but its effect on selection accuracy depends on retrieval coverage and routing
errors as the available library grows.

\section{Problem and related work}
A reusable skill can package task instructions in a \texttt{SKILL.md} document.
Showing every skill description makes the discovery prompt grow with the library.
Retrieving a shortlist controls this input, but can remove a necessary skill
before the model sees it. Routing to one category narrows the search further,
while introducing another possible point of failure. We ask when this trade-off
is useful, without assuming that any strategy must win.

SkillRet supplies a large skill-retrieval benchmark and highlights the difficulty
of finding relevant skills in noisy requests \citep{kang2026skillret}. Skill
Retrieval Augmentation separates retrieval, skill incorporation and downstream
task execution \citep{su2026skill}; our project studies the first two boundaries
through retrieval and metadata-only selection, not task completion.
Work on single-agent skill libraries motivates examining semantic confusion and
hierarchical routing as libraries expand \citep{li2026single}. Skill Shadowing
separates wrong skill selection from execution degradation caused by additional
context \citep{song2026shadowing}. Graph-of-Skills investigates dependencies and
structured retrieval \citep{liu2026graph}; dependency graphs and skill chains are
outside our scope. Our contribution will be a reproducible, bounded comparison
of simple discovery strategies using real native skill middleware, with failures
and costs reported alongside accuracy.

\section{Experimental design}
\paragraph{Data and library construction.}
We use the pinned SkillRet test snapshot with 6,006 skills, 4,392 synthetic queries
and 7,187 relevance judgments. Its six-way \texttt{major} taxonomy supports the
router. For each query, the sampler retains all annotated relevant skills and
adds seeded random distractors to form nested libraries of $N=50,250,1000$.
Membership and display order are shared across strategies. This is a conditional
distractor-scaling study, not the official full-corpus benchmark or a production
catalogue. Labels are available only to the sampler and evaluator; the model sees
canonical query text and candidate metadata.

\paragraph{Four methods, one selector.}
\textbf{Listing} exposes all $N$ descriptions. \textbf{BM25} retrieves up to five
candidates using keyword matching. \textbf{Vector} retrieves five by cosine
similarity using the pinned CPU BGE-small-en-v1.5 encoder.
\textbf{Domain-vector} asks the same language model to predict one of six major
categories, then searches within that category. An empty category is recorded as
a routing failure, without global fallback. Each arm uses identical original
name/description text, capped once at 1,024 characters and 128 BGE tokens.

\paragraph{Native integration and controls.}
Candidates enter the real LangChain Deep Agents \texttt{SkillsMiddleware} through
an isolated in-memory \texttt{StateBackend}. Main selection allows one
\texttt{select\_skills} call returning zero to five original IDs; filesystem,
shell and subagent tools are excluded. A separate activation check permits
bounded native reads, never execution of downloaded instructions. Temperature,
output limit, prompt, provider, embedding revision and candidate order are fixed
across methods. We propose DeepSeek V4.1 Flash on its pinned provider as the
primary model because its pilot integration completed most cases, not because
the pilot established superior accuracy. Availability must pass a new canary.

\begin{figure}[t]
\centering
\fbox{\parbox{0.93\linewidth}{\centering
Canonical request + seeded skill library\\[3pt]
$\downarrow$ listing / BM25 / vector / domain-vector $\downarrow$\\[3pt]
Candidate metadata $\rightarrow$ native skill middleware $\rightarrow$ selected IDs\\[3pt]
$\downarrow$ evaluator-only labels, saved failures, tokens and costs}}
\caption{The same selection boundary follows every discovery method. Skill bodies
are read only in the separate activation check.}
\label{fig:flow}
\end{figure}

\section{Executed pilot and feasibility}
% These are observed snapshot values. Do not replace them with a planned result.
The real baseline ran before API spending. Table~\ref{tab:retrieval} separates
full-document retrieval from the common-metadata representation. The former is
a 200-query partial-protocol check, not reproduction of a published full-test
score. All rows use the same 200 consumed pilot queries and all 6,006 test skills.

\begin{table}[h]
\centering
\begin{tabular}{llr}
\toprule
Retriever & Document representation & Recall@5 \\
\midrule
BM25 & Full name, description and body & 51.67\% \\
BM25 & Capped name and description & 60.58\% \\
BGE-small & Same capped metadata & 49.25\% \\
\bottomrule
\end{tabular}
\caption{Measured retrieval baselines. Input representation changes between the
first and second rows, so that difference is not a method-only comparison.}
\label{tab:retrieval}
\end{table}

The live follow-up compared five fixed OpenRouter model/provider configurations
on three requests, three library sizes and four methods: 180 planned cases,
138 attempted, 121 valid selections and 24 exact reference sets. Of these cases,
53 reused compatible earlier records. At $N=1000$, none of 35 valid selections
was exact, although 22 included at least one annotated skill. Qwen 27B rate limits,
MiMo overload and invalid IDs reduced coverage. These unequal, tiny denominators
do not support a general model ranking or a reliable scaling law.

All five revised integration canaries passed. Separate activation recorded 16
actual native reads across 25 cases, with 15 completed workflows. Total recorded
project spending was US\$0.3865, or an estimated CAD\$0.5797 at the declared 1.50
planning factor. This includes earlier stages, not just the scaling comparison.
The reproducible handoff contains two notebooks, 54 passing tests, exact source
pins and sanitized traces. Code and measured evidence are available in the
\href{https://github.com/speccy88/mia5150/tree/b9c9df2/skillret-pilot}{versioned pilot folder}.

\section{Evaluation and next-stage plan}
\paragraph{Metrics and interpretation.}
Let $G_q$ be the annotated skill set, $C_q$ the retrieved candidates and $S_q$ the
valid selected set. We will report candidate recall $|C_q\cap G_q|/|G_q|$ separately
from selection precision, recall and F1. Exact-set match requires $S_q=G_q$.
Every table will show planned, attempted and valid counts. Conditional accuracy
will be accompanied by exact successes per planned case, so provider exclusions
and malformed outputs remain visible rather than disappearing from the denominator.
Token accounting includes input and output, plus routing in every logical
domain-vector workflow even when a physical route is reused. Actual USD charges,
estimated CAD, latency and failed-call reservations will be reported separately.

\paragraph{Sampling and analysis.}
The 200 exploratory queries remain consumed; 4,192 disjoint final IDs are reserved.
The next bounded batch will preregister 20 previously unused queries, including
single- and multiple-positive cases, with seed 42 and the three fixed library
sizes. It requires at most 240 selectors and 20 routers before exclusions, plus
canary/activation calls. Additional seeds 43 and 44 are planned only after an
explicit budget review. We will compare methods on paired queries, report
uncertainty over queries when sample size permits, and inspect missing candidates,
wrong categories and extra selections separately. A more independent domain-vector
arm should replace the pilot's earlier-output gate in a separately versioned
protocol; historical gated results will not be pooled with that new arm.

\paragraph{Resources, risks and backup.}
The existing harness enforces US\$3 and CAD\$5 ceilings, a persistent 300-request
cap, per-call reserves and no automatic retries. Its recorded 236 attempts leave
insufficient capacity for the proposed next batch: a separately authorized study
budget is required, and the existing ledger will not be reset. No new paid calls
are authorized by this proposal. Provider availability and context limits may
reduce coverage; we will retain exclusions and use one stable endpoint. Synthetic
queries, imperfect relevance labels, random distractors and seven truncated BGE
query embeddings in the pilot limit generalization. If live coverage remains
unreliable, the backup is full-corpus BM25 versus BGE on matched metadata, followed
by a smaller $N=50$ native selection study. No training or GPU is required.

\begin{table}[h]
\centering
\small
\begin{tabular}{p{0.12\linewidth}p{0.55\linewidth}p{0.23\linewidth}}
\toprule
Weeks & Deliverable & Proposed lead \\
\midrule
5--7 & Freeze queries, protocol and provider; check reproducibility & Anonymous \\
8--10 & Run authorized batches; audit labels, costs and exclusions & Anonymous \\
11 & Paired analysis, error taxonomy and figures & Anonymous \\
12--13 & Presentation, report and reproduction review & All three \\
\bottomrule
\end{tabular}
\caption{Proposed responsibilities, to be confirmed by the team.}
\label{tab:timeline}
\end{table}

\noindent\textbf{Reproducibility and disclosure.}
We will retain source/configuration hashes, query reservations and failures; no
labels enter model state and no skill instructions execute. AI tools supported
code preparation and proposal drafting. Team review of the methods, references,
roles and final claims is required before submission.

% References do not count toward the course's 2--3 page proposal body.
\clearpage
\bibliographystyle{plainnat}
\bibliography{references}
\end{document}

